\chapter{April}

\section[April 6th]{April 6\textsuperscript{th}}

Today's book is called ``The Rainbow Fish'', by Marcus Pfister. It is about a prideful fish with shiny and colored scales that thought he was more beautiful than the others and had no friends. The Rainbow Fish started to share its colored scales with other fishes until only one remained with him. With that, he made many friends. It is a nice and quick story for children. This time it took less than 15 minutes; today it served as a decent icebreaker to start the class. I didn't find the book to be that deep. Of course, it is important to share the belongings we have plentiful of and care about others.

\textbf{Advice for quotations:} Double quotations are used for direct citations to a text with the reference coming right after. Single quotations are not an exact quotation, but indicate the idea is not original and comes from another text, or represent mere linguistic expressions.

For today's class, we should have read the third chapter of Tomitch's book about what a text is. It is interesting in the sense that it says a text is a unit of language in use with meaning. A stop sign, a no-smoking sign, or Post-it notes are all texts. They are ``instances of language in use that communicate something to the reader'' as the slides state.

The professor gave an example of the expression ``robot ahead''. An important author in English language studies saw a sign similar to a stop sign with that text in South Africa; there he learned that they call traffic lights robots.

98\% of Brazilians don't think they can speak English. This reminds me of the bad experience I had with English teaching in my school, where the teacher would spend most of the time showing movies and doing really easy tests to the point where most of my colleagues simply didn't learn the language. Most people also don't see the utility of it, ignoring the benefits for the brain, and the world it opens beyond travel, such as in research, studies, job opportunities, and even online entertainment.

The professor made an argument about how Blumenau and the region cannot call themselves a small Germany outside of Germany; we are Brazilians. On this topic, I always remember this video about Blumenau in 1970: \url{https://www.youtube.com/watch?v=BYiHkU7KxgU}. In it, Blumenau people are called ``brasileiríssimos descendentes dos colonizadores (europeus)''. It sounds weird to me when it is said that we shouldn't reference Germany when describing the people from Blumenau and the region, and simply call us Brazilians.

The main argument of the teacher is that the German language is not as prominent in the city nowadays. As I understand it, language is just one of many manifestations a culture can have. It is a bit difficult to measure the influence of German culture in the region by looking at the language alone, because it was banned by the dictator Getúlio Vargas. This means that the language wasn't abandoned due to it losing importance; instead, it was an authoritarian mandate from a dictator that artificially suppressed the language. Due to this external interference, the German language ceased to be an accurate measurement of how German the people from Blumenau really are.

When I talk about my upbringing with people from other states (and even other regions of Santa Catarina like Itajaí, Florianópolis, or Balneário Camboriú), they consider it to be really different. My grandparents (whom I still call ``oma'' and ``opa'') prefer to speak German with my relatives. The culinary manifestations of culture with the ``cuca'', the architecture, the emphasis on keeping beautiful and organized gardens, making jams (attached to the verb ``schmiar'', similar to ``smear'' in English), the different kinds of breads (manioc bread, potato bread, and so on), the relationship with Germany in our universities (e.g., Letras Alemão), companies with ties to Germany (T-Systems and EB-ISCO), and even the German Consulate. Of course, there are also many differences from modern-day Germany, but the influence still remains strong in my view. Even though we are still in Brazilian territory and shouldn't be considered outside of it, we also aren't only Brazilian in our identity.

``Never be ashamed that we are Brazilian''. We tend to judge ourselves and others based on English.

\begin{description}
    \item[Cohesion $\rightarrow$ Micro-level:] Concerns the grammatical relationships within a text, specifically how words and ideas connect. Important elements are pronouns and conjunctions. It refers to the syntax. ``The ways words connect to make sense in a sequence'' (Words by Professor Bailer).
    \item[Coherence $\rightarrow$ Macro-level:] Goes beyond the surface of the text; it refers to the ideas and concepts outside of the constraints of the words themselves. The real world and the general idea of the text (i.e., Global Meaning) should be identifiable from the relations of the words.
\end{description}


A logical example of an argument is:
\begin{quote}
    Socrates is a man.\\
    All men are mortal.\\
    Therefore, Socrates is mortal.
\end{quote}

An incoherent alternative would be:
\begin{quote}
    Socrates is a man.\\
    Men are rational.\\
    Therefore, Socrates is mortal.
\end{quote}

It is a \textit{non sequitur}; the sentences are related (either to man or to Socrates), but the argument makes no sense.

Every text must also have the following attributes:
\begin{description}
    \item[Intentionality:] The author must have an intention to express an idea through the text.
    \item[Acceptability:] Refers to the receiver of the text. The text must have some sort of relevance to the receiver who needs to ``uncover the meaning'' intended by the writer.
    \item[Informativity:] The text must be able to inform an idea that isn't completely unrelated to the receiver's understanding, but should also present new information so as not to cause boredom.
    \item[Situationality:] Presents the context in which this text is relevant to be used (e.g., a stop sign by the side of the road vs. in front of a teenager's door).
    \item[Intertextuality:] The relation of the current text with other texts; this is a relation of dependencies on the information in other texts.
\end{description}

\section[April 13th]{April 13\textsuperscript{th}}

Today's story session also wasn't too long. It is called ``Not a Box'', a book about a bunny that uses imagination to play with a box, he sees the box as many things such as a building and an airplane. The lesson is about the importance of having a good imagination and thinking ``outside the box''. I don't think anyone would undermine the importance of creativity when studying and doing science, of course, within the boundaries of common sense so we don't waste time with baseless crazy ideas.

Next, we talked about procrastination. I relate to it quite a bit: daydreaming, napping, doomscrolling, gaming, online shopping, and so on. For sure, it is one of the problems that afflict a huge percentage of society. The first thing that comes to my mind is the ``student syndrome'' (for more information, you can refer to the article by Dr. Deysolong at: \url{https://www.researchgate.net/publication/370731697_STUDENT_SYNDROME}). Basically, we wait until the last possible moment to start a project. It is difficult to say something about procrastination that hasn't been said before to exhaustion; nevertheless, here is an animation that conveys the idea well in a highly relatable way: \url{https://www.youtube.com/watch?v=F8zqVO8UqTM}.

Professor Bailer suggests creating an organized agenda for the day and trying our best to stick to it. To me, it looks like this idea can fall into a trap. After all, preparing the agenda and planning our life is in itself a project and work; even if it is done daily, weekly, or monthly, I see myself procrastinating the schedule creation itself. For now, I am still managing everything in my head; I did that throughout elementary, middle, and high school, as well as during my CS major.

Now we entered the official discussion of the text. It was about chapter 4 of Tomitch's book, mainly differentiating essays and research papers. An essay can be similar to a composition we write in high school, while extending into longer and more complex formats. It follows the structure with an introduction, development, and conclusion. It is acceptable to include more of one's own opinions.

In the research paper, there is more documentation and research itself; the writer needs to do a broad study on an area and map everything to the topic being discussed. Then it creates a ``hole'' in the theory, in which the paper itself can fit and contribute to the production in that area. It contains an original research idea and the key findings of this research.

I like the distinction between different kinds of essays and papers. Especially the reflection about how in high school we write an essay for only one reader (the teacher), who already knows everything about the topic. Later, we start to write both essays and papers for the whole community, sharing new information they do not know yet. This is quite a big difference in responsibility and in the way of thinking when structuring the text.

In the research paper, we have the hourglass format: the introduction starts general until you dig your ``hole'' and work on it, then you have the discussion where the findings are discussed from a big-picture perspective, showing how a piece fits into the puzzle that is the area being studied.

The next topic is conjunctions. I am familiar with this term from my logic classes. Conjunction holds the value of the clause ``and'' ($\land$ symbol). The opposite is disjunction, which means the value ``or'' ($\lor$ symbol). \sloppy Here is a more comprehensive text about logic with conjunction and disjunction from Stanford University: \url{http://intrologic.stanford.edu/chapters/chapter_02.html}.\fussy

Of course, it is a bit different in linguistics. In text writing, conjunctions are words used to connect different sentences. Conjunctions have different relations like addition, opposition, justification, and so on. It looks like ``and'' is the most basic conjunction, a base. ``X and Y'' sounds raw. If we have ``X but Y'', logically speaking, you are still stating both X and Y while indicating to the reader a relation of opposition. With ``X because Y'', once again, the logical value remains (the speaker is affirming both X and Y), while from a reader's perspective, those words provide more information.

There are coordinating conjunctions; they join words, phrases, and clauses:
\begin{itemize}
    \item \textbf{F} = For
    \item \textbf{A} = And
    \item \textbf{N} = Nor
    \item \textbf{B} = But
    \item \textbf{O} = Or
    \item \textbf{Y} = Yet
    \item \textbf{S} = So
\end{itemize}

A big difference is that in this context, ``or'' is a conjunction. In logic, it fits squarely in the disjunction category.

Subordinate conjunctions join independent and dependent clauses. They convey a more complex relationship between the clauses:
\begin{itemize}
    \item I can stay out \textit{(independent)} until the clock strikes twelve \textit{(dependent)}.
    \item Before he leaves \textit{(dependent)}, make sure his room is clean \textit{(independent)}.
    \item I drank a glass of water \textit{(independent)} because I was thirsty \textit{(dependent)}.
    \item Because I was thirsty \textit{(dependent)}, I drank a glass of water \textit{(independent)}.
\end{itemize}


\section[April 27th]{April 27\textsuperscript{th}}

For today, our homework was to read chapter 5 of Tomitch's book on Academic Writing. The focus was on how to summarize academic texts. The thought that we are already biologically capable of summarizing events, like movies, past events, and so on, was particularly interesting to me. Summarizing complex ideas from an academic text is more mentally challenging, because it lacks the familiar narrative structure, often presenting a less straightforward organization. The actual tips and thoughts about how to summarize, alongside the realization that we remain authors ourselves when using other authors to write our texts, make a lot of sense. The perspective and choices we make about which authors to use and how to harmonize their works are really interesting.

Today's children's book is ``How to Catch a Star'' by Oliver Jeffers. It was about a child that liked stars and waited the whole day for a star until it appeared at night. Then he wanted to reach the star; it turned out he couldn't. So he waited by the seashore and caught what looked like a starfish. It seems that the story is about determination and how people should focus on their objectives without giving up. Today it took quite a while to finish the story again; the story time only ended at about 7:20 pm.

The professor brought some reflections with the help of ChatGPT, talking about how the imagination of the child can inspire us to be more persistent, creative, and willing to try different points of view to reach or adapt our objectives. She asked if the book helps us with that, referring to our professions (for me in CS, for most of the class, teachers). Honestly speaking, I don't think so, not specifically at least. If the book can help only after intensive reflective and metaphorical thinking, then anything can. We can reflect on almost anything and reach meaningful conclusions. This approach relies entirely on our own thought process rather than providing an in-depth framework for seeing life; thus, the book itself has little merit in this so-called ``career help''. Perhaps it is more useful as a means of making those life lessons interesting to children, so in that niche, I get its usefulness.

Continuing with the analysis of the 5\textsuperscript{th} chapter from Tomitch, a new concept learned is ``salience'', the notion of how important a piece of information is in relation to the current context we are in.

We need to start thinking about the summary before starting the reading. We need to reflect upon the type of summary we want to make and what use it will have. The analysis must be done in terms of introduction, methods/key points, results, and discussion (for articles), and introduction, development, and conclusion (for essays).

In English, we have a distinction between an abstract and a summary. The former is academic, while the latter might not be; this distinction doesn't exist in Portuguese.

We always need to be careful when delivering a text. After all, we send it by itself and lack the opportunity to explain what we meant when writing. Therefore, it must always be revised.

We also checked the concepts of microlevel and macrolevel. The first refers to the actual sentences of the text itself, while the macrolevel is a ``cleaner'' version, where we remove trivial information, delete redundant information, superordinate lists and actions, and select or invent a topic sentence.

Another important concept is paraphrasing. A tip is to read the original text, understand the idea, and then try to recite the main ideas to yourself without checking the original material. Then, we write the paraphrase that we thought about, and finally, we can check the paragraph against the original to see if it is acceptable. If we cannot succeed in paraphrasing, we should opt for direct citation.

Summaries can be used to compose part of our own text.

I like Tomitch's concept of ``sieving and sowing''; we sieve the academic texts to select the most relevant ones, and then sow what we selected to make them harmonize and contribute towards the main idea we are proposing.

Three steps to summarize:
\begin{description}
    \item[R: Read.] \hfill
    \begin{itemize}
        \item Read to get familiar
        \item Read out loud
        \item Read as a believer
        \item Read as a doubter
    \end{itemize}
    \item[A: Annotate.] \hfill
    \begin{itemize}
        \item Write in margins
        \item Underline important parts
        \item Use sticky notes
        \item Circle words you don't understand
    \end{itemize}
    \item[T: Think.] \hfill
    \begin{itemize}
        \item What is the purpose of the author?
        \item What are the facts in the text?
        \item What assumptions are made?
    \end{itemize}
\end{description}

The summary might be about one-third of the original text (as a rule of thumb).

This class has been very useful in providing a better idea of how to start with summaries; I didn't think there was such a strict method to write them. Usually, I just summarize blindly when I need to, trying to list the ideas that I can remember in a structured manner. In reality, things go way beyond that, especially regarding the techniques to avoid accidental plagiarism caused by repeating the same sentences from the original with synonyms.

The professor talked about writing in law; she said something along the lines of, ``If we have so many crimes inside the law, that means that the law isn't well written enough; after all, it can be bended to commit crimes without consequences.'' I don't necessarily agree with that. The way the law is set up in Brazil follows Post-Positivism philosophy. In summary, it contrasts Natural Law and positivism. Positivism is the idea that we should strictly follow the law exactly as it is written; because the Nazis used this argument to justify their actions as simply ``following orders'', the legal system decided to adopt an alternative approach.

The idea is that we should go beyond what is written in the text and ``feel'' the way society is advancing, interpreting the law through the lenses of the morals that guide the constitution. The literal written word of the law can be ignored if it contradicts what is understood as the ``spirit of the constitution'' and the morality we should aim for. The evolutions in the understanding of Human Rights should be followed even if new laws reflecting them aren't explicitly created by the legislative power. For more information, there is a nice text by the former STF judge Luís Roberto Barroso commenting on it here: \url{https://periodicos.fgv.br/rda/article/view/43618/44695} (parts 1 and 2 should be enough to get the gist of the interpretation of the law in Brazil currently). Following this logic, properly written texts can still be reinterpreted beyond their literal wording, changing the most intuitive meaning according to the general values that guide the Brazilian constitution, as determined by the judges.

We also took a look at ``corpus linguistics'', which is really interesting: a large collection of texts that can be analyzed to check the frequency of terms and get ideas about how the language is structured. Such a massive amount of texts in both Portuguese and English appeals to someone from Computer Science like me; I instantly think about training my own AI on the language to see if I can talk with my creation and observe how the corpus will teach it.